\subsection{The free Cartesian equation}

Put \(\mathcal T=\partial_\tau-y\cdot\nabla/2\),
\(L=\mathcal T-\nu\Delta-1/2\), and let \(Y_r\) denote
\(L^r((\mathbb R/P\mathbb Z)\times\mathbb R^3)\), with normalized
phase measure. Its maximal distributional graph is
\[
 D_r=\{w\in Y_r:\nu\Delta w,\mathcal Tw\in Y_r\},\qquad
 \|w\|_{D_r}=\|w\|_{Y_r}+\|\nu\Delta w\|_{Y_r}
                                      +\|\mathcal Tw\|_{Y_r}.
\]
The derivatives in this definition are Cartesian; there is no axis
condition. Write \(\gamma_r=(3/r-1)/2\), \(a_r=\operatorname{Re}\mu+\gamma_r\).

\begin{lemma}\label{bg:common-body-graph}
For \(1<r<3\), \(a_r>0\), \(\nu>0\), and \(0<P\le1\),
the componentwise operator \(\mu+L:D_r\to Y_r\) is bijective and
\begin{equation}\label{bg:body-graph}
 \|w\|_{D_r}+\sqrt\nu\|\nabla w\|_{Y_r}
 \le C_r[1+(1+|\mu|)/a_r]\|f\|_{Y_r},
 \qquad (\mu+L)w=f.
\end{equation}
Its projected inverse is obtained by applying Cartesian Leray projection
to \(f\). At \(r=2\) the adjoint has the same maximal graph and the
same bound. All these estimates are independent of \(P\).
\end{lemma}

\begin{proof}
The heat-dilation formula gives a distributional inverse
\[
 S(t)h(y)=e^{t/2}[e^{\nu(e^t-1)\Delta}h](e^{t/2}y),\qquad
 \mathcal R_\mu f(\tau)=\int_0^\infty e^{-\mu t}S(t)f(\tau-t)\,dt.
\]
Its integrand has \(Y_r\) norm at most
\(e^{-a_rt}\|f\|_{Y_r}\); hence \(\|\mathcal R_\mu\|\le a_r^{-1}\).
The local heat estimate needed here is
\begin{equation}\label{free:heat-estimate}
 \|\nu D^2u\|_r+\|u_t\|_r\le C_{d,r}\|F\|_r,\qquad
 u_t-\nu\Delta u=F,
\end{equation}
for compactly time-supported distributional solutions \(u,F\in
L^r(\mathbb R\times\mathbb R^d)\). To justify its precise forcing
class, scale to \(\nu=1\). The multiplier
\(-\xi_i\xi_j/(i\omega+|\xi|^2)\) has \(L^2\) norm at most one.
Its kernel away from the origin is
\(K(t,x)=1_{t>0}\partial_{ij}H_t(x)\). For
\(\rho(t,x)=\max\{|x|,\sqrt{|t|}\}\), Gaussian differentiation gives
\[
 |K|\le C\rho^{-d-2},\quad
 |\nabla_xK|\le C\rho^{-d-3},\quad
 |\partial_tK|\le C\rho^{-d-4}.
\]
The kernel extends smoothly by zero through \(t=0\), away from
\(x=0\). Integrating these derivatives in parabolic annuli gives
\[
 \int_{\rho(z)>4\rho(h)}|K(z-h)-K(z)|\,dz\le C_d.
\]
The spatial and temporal differences are bounded, respectively, by
\[
 C|h_x|\int_{4\rho(h)}^\infty R^{-2}dR,\qquad
 C|h_t|\int_{4\rho(h)}^\infty R^{-3}dR,
\]
which are uniformly bounded.
For completeness, decompose an \(L^1\cap L^2\) function at height
\(\alpha\) using maximal dyadic rectangles with spatial side \(2^k\)
and time side \(2^{2k}\). The good part has
\(\|g\|_2^2\le C_d\alpha\|f\|_1\). The bad parts \(b_Q\) have
zero integral, total \(L^1\) norm at most \(2\|f\|_1\), and their
rectangles have total measure at most \(\alpha^{-1}\|f\|_1\).
Outside fixed enlargements of these rectangles, cancellation and
the displayed kernel integral bound the \(L^1\) norm of each
\(Tb_Q\) by \(C_d\|b_Q\|_1\). The \(L^2\) estimate for the good
part and Chebyshev's inequality for the bad part prove weak type
\((1,1)\). Integrating this distribution bound for the portions
\(|f|>\alpha\) and \(|f|\le\alpha\), together with the \(L^2\)
bound, gives \(L^r\) boundedness for \(1<r<2\). Duality gives the
remaining exponents since the reflected kernel has the same bounds.
Mollify the compact-time equation. The mollified solution vanishes
in the distant past, so its causal heat formula identifies its
second derivatives with these bounded operators applied to the
mollified \(F\). Their \(L^r\) limit equals the distributional
derivatives of \(u\); the equation gives \(u_t\). This proves
\eqref{free:heat-estimate}, with no viscosity loss.

On \(|t|<2\), set
\(v(t,x)=e^{-t/2}w(\tau_0+t,e^{-t/2}x)\).
Then \(v_t-\nu e^t\Delta v\) is the same pullback of \(f-\mu w\).
A fixed cutoff equal to one on \(|t|<1\), followed by
\(\theta=e^t\), reduces this equation to the preceding heat estimate.
All Jacobians on the window are bounded. Consequently
\[
 \int_{\tau_0-1}^{\tau_0+1}\|\nu\Delta w\|_r^r\,d\tau
 \le C_r\int_{\tau_0-2}^{\tau_0+2}
       [\|f\|_r^r+(1+|\mu|)^r\|w\|_r^r]\,d\tau.
\]
For periodic \(h\),
\(\int_0^P\int_Ih(\tau_0+t)\,dt\,d\tau_0
=|I|\int_0^Ph\).
Averaging the last inequality over the centers proves the diffusion
bound with no inverse-period factor. The equation gives the drift
bound; spatial elliptic interpolation gives the gradient bound.
The same reasoning applies to any distributional \(Y_r\) solution.

The adjoint semigroup is
\(S(t)^*h=e^{-t}[e^{\nu(1-e^{-t})\Delta}h](e^{-t/2}y)\).
For a smooth compact spatial test, its adjoint resolvent has itself and
its first two spatial derivatives in \(Y_{r'}\), by the gap and the
extra factor \(e^{-t/2}\) per derivative. Test a homogeneous solution
against this resolvent times a radial cutoff at radius \(N\).
The drift cutoff is bounded, and the diffusion cutoffs cost
\(N^{-1},N^{-2}\); all terms are bounded by the \(Y_r\) tail of the
solution times fixed \(Y_{r'}\) norms and tend to zero. Thus the
solution vanishes. The inverse therefore has precisely the displayed
maximal distributional domain. To exhibit a core directly, approximate
the right side by smooth periodic Schwartz functions and truncate the
inverse integral at \(M\). Such a truncated integral is Schwartz and
its equation has right side
\[
 f-e^{-\mu M}S(M)f(\tau-M).
\]
The remainder tends to zero in \(Y_r\); the graph estimate gives
graph convergence of the truncations. Thus the smooth periodic
Schwartz class is a core. Radial cutoffs approximate it in this graph,
since their drift commutator is bounded on an escaping annulus and
their diffusion commutators tend to zero. At \(r=2\), reverse
the phase window and use \(\theta=1-e^{-t}\) for
\(L^*=-\mathcal T-\nu\Delta+1\). Testing against the compact core
identifies the distributional equation of every vector in the actual
adjoint domain; this reversed-window estimate places it in \(D_2\).
Conversely, core integration by parts puts every element of \(D_2\)
in the adjoint domain. Its resolvent norm is at most \(a_2^{-1}\)
by duality, so the same graph bound follows. Leray commutes with
heat, dilation and phase translation, so projection preserves all
identities. Restriction to the real axisymmetric no-swirl subspace is
also legitimate since its symmetry projections commute with these maps.
On the Hilbert core, direct integration by parts gives
\[
 \operatorname{Re}\langle(\mu+L)w,w\rangle
       =(\operatorname{Re}\mu+1/4)\|w\|_2^2
                                      +\nu\|\nabla w\|_2^2.
\]
The adjoint has the same real part. For a tensor source
\(\mathbb P\operatorname{div}T\), its right side is at most
\(\|T\|_2\|\nabla w\|_2\). Completing the square gives
\(\|\nabla w\|_2\le\nu^{-1}\|T\|_2\) and
\(\|w\|_2\le[2\sqrt{\nu(\operatorname{Re}\mu+1/4)}]^{-1}\|T\|_2\).
Approximation of tensors completes this weak-stress map and the
inequality without assigning a full graph to its rough outputs.
\end{proof}

\begin{lemma}\label{bg:graph-embeddings}
Every \(w\in D_r\) has a representative in
\(C(\mathbb R/P\mathbb Z;L^r)\), and
\begin{equation}\label{bg:phase-embedding}
 \sup_\tau\|w(\tau)\|_r
 \le(1+3P/(2r))\|w\|_{Y_r}+P\|\mathcal Tw\|_{Y_r}.
\end{equation}
For \(r>5/2\),
\begin{equation}\label{bg:supremum-embedding}
 \|w\|_{L^\infty_{\tau,y}}
       \le C_r\nu^{-3/(2r)}\|w\|_{D_r}.
\end{equation}
\end{lemma}

\begin{proof}
On a fixed phase interval the distributional chain rule gives
\(V(t,x)=w(\tau_0+t,e^{-t/2}x)\in W^{1,r}_tL^r_x\), with
\(V_t\) the pullback of \(\mathcal Tw\). Undoing the strongly
continuous dilation gives the continuous phase representative;
distributional periodicity matches its endpoints. For
\(f(\tau)=\|w(\tau)\|_r\), the norm chain rule gives
\(\lvert f'+3f/(2r)\rvert\le\|\mathcal Tw(\tau)\|_r\).
Start at a phase with \(f\le\|w\|_{Y_r}\), integrate for at most
one period, and apply H\"older to obtain \eqref{bg:phase-embedding}.

For the supremum fix \((\tau_0,y_0)\) and put
\(W(t,z)=w(\tau_0+t,e^{-t/2}(y_0+\sqrt\nu z))\) on
\(|t|<1/2\). Then \(W_t=(\mathcal Tw)\circ\),
\(\Delta_zW=\nu e^{-t}(\Delta_yw)\circ\), and
\(dz=\nu^{-3/2}e^{3t/2}dy\). The integral of a nonnegative
\(P\)-periodic function over a unit interval is at most
\((1+P)\) times its phase average. Thus
\[
 \|W\|_{L^r}+\|W_t\|_{L^r}+\|\Delta_zW\|_{L^r}
       \le C_r\nu^{-3/(2r)}\|w\|_{D_r}.
\]
Choose a fixed time cutoff \(\chi=1\) near zero. The right side
\(F=(\partial_t-\Delta_z+1)(\chi W)\), after zero extension,
has this same bound. Its causal inverse kernel
\(k(t,z)=1_{t>0}e^{-t}H_t(z)\) satisfies
\[
 \|k\|_{L^{r'}}^{r'}
   =C_r\int_0^\infty e^{-r't}t^{-3(r'-1)/2}\,dt<\infty
       \quad\Longleftrightarrow\quad r>5/2.
\]
The identity \(\chi W=k*F\) follows distributionally; uniqueness
follows from the nonvanishing symbol \(i\omega+|\xi|^2+1\).
H\"older bounds its value at \((0,0)\), proving
\eqref{bg:supremum-embedding}. Translation continuity of \(k\) in
\(L^{r'}\) supplies a continuous local representative. No external
Sobolev embedding or inverse-period factor is used.
\end{proof}

\subsection{Tensor forces}

Write \(R_\mu=(\mu+L)^{-1}\mathbb P\). For tensor inputs its
divergence is interpreted by completion from smooth tensors, rather
than as an ordinary \(Y_2\) body force. Let
\[
 B=\mathcal T-\tfrac34,\qquad A=-\nu\Delta.
\]
The group \(u(\tau,y)\mapsto e^{-3t/4}u(\tau+t,e^{-t/2}y)\) is
unitary on \(Y_2\) and has generator \(B\); hence \(B\) is skew-adjoint.
The operator \(A\) is nonnegative self-adjoint.

\begin{lemma}\label{free:weak-stress}
For \(a=\operatorname{Re}\mu+1/4>0\), the map
\(S_\mu=R_\mu\operatorname{div}\) extends uniquely to
\(Y_2(\operatorname{Sym}_3)\), with
\begin{equation}\label{free:stress-energy}
 \|S_\mu T\|_{Y_2}\le\frac{\|T\|_{Y_2}}{2\sqrt{a\nu}},\qquad
 \|\nabla S_\mu T\|_{Y_2}\le\nu^{-1}\|T\|_{Y_2}.
\end{equation}
Uniformly on compact subsets of this half-plane,
\begin{equation}\label{free:half-drift}
 \|(1+|B|)^{1/2}S_\mu T\|_{Y_2}
       \le C_\mu\nu^{-1/2}\|T\|_{Y_2}.
\end{equation}
The weak solution is unique among distributional \(Y_2\) solutions.
Bounded measurable tensor partitions are allowed before divergence.
The body and tensor inverses depend
holomorphically on \(\mu\) in their respective half-planes.
\end{lemma}

\begin{proof}
The energy identity and completion proved above give
\eqref{free:stress-energy}. Differences of smooth tensors give strong
convergence of both velocities and first gradients, so the equation
and energy identity pass to the limit. Distributional uniqueness
follows by subtracting two solutions and applying the homogeneous
uniqueness argument in Lemma~\ref{bg:common-body-graph}.
In particular a bounded measurable multiplier may be applied to a
tensor before this completion; the distributional divergence includes
all interface terms.

Put \(C=1+|B|\), \(D=1+A\), and let \(\mathcal R_\mu\) be the
inverse of \(\mu+L\), acting componentwise. Its direct graph bounds
\(C\mathcal R_\mu\); the actual adjoint graph bounds
\(D\mathcal R_\mu^*\), and duality bounds the extension of
\(\mathcal R_\mu D\). The imaginary powers of \(C,D\) are unitary.
Insert the respective spectral cutoffs on the two sides of
\[
 C^z\mathcal R_\mu D^{1-z},\qquad
                 0\le\operatorname{Re}z\le1.
\]
At the two boundary lines the preceding graph estimates give
uniform bounds. The scalar three-lines theorem applied to pairings,
followed by density and closedness, yields
\(\|C^{1/2}\mathcal R_\mu D^{1/2}\|\le C_\mu\).
No commutation between \(A\) and \(B\) is asserted. Spatial Fourier
transformation gives
\(\operatorname{div}=\nu^{-1/2}A^{1/2}\mathcal Q\), where
\(\mathcal Q\) is a bounded tensor-to-vector combination of Riesz
transforms. Since \(D^{-1/2}A^{1/2}\) is bounded and Leray commutes
with \(A,B\), the last estimate proves \eqref{free:half-drift}.

The inverse of \(\mu+L\) is holomorphic by its convergent semigroup integral.
For the stress inverse, distributional uniqueness gives
\[
 S_\mu-S_\zeta=(\zeta-\mu)R_\mu S_\zeta.
\]
The body and stress bounds make this the resolvent identity in
operator norm, proving the asserted holomorphy.
\end{proof}

At \(r=11/4\) the body gap at \(\mu=0\) is \(1/22\), and at
\(r=23/8\) it is \(1/46\). In particular all shifts in the
background and spectral constructions lie in the stipulated ranges.
The graph embeddings also give, at \(r=11/4\),
\[
 \|\mathbb P[(u\cdot\nabla)v]\|_{Y_r}
 \le C\nu^{-23/22}\|u\|_{D_r}\|v\|_{D_r}.
\]
For any unitary velocity and tensor lifts \(J_v:H\to Y_2\) and
\(J_T:H_T\to Y_2\), the transported stress map is exactly
\(J_v^{-1}S_\mu J_T:H_T\to H\), with the same energy bounds.
This formula transports the completed tensor inverse; it does not
assign an ordinary body-force norm to \(\operatorname{div}J_TT\).
